Theorem 1 by removing noisy causal patterns and learning
the optimal baseline value, so as to boost the sparsity of
causal effects. Besides, we also discovered that adversarial
training [46] can make the DNN encode much more sparse
causal effects.
First, boosting conciseness by learning the optimal
baseline value. In fact, the sparsity of causal patterns does
not only depend on the DNN itself, but it is also deter-
mined by the choice of baseline values in Eq. (4). Specif-
ically, input variables are masked by their baseline values
r = [r1,r 2,...,r n] to represent their absence states in the
computation of causal effects. Thus, wΩ can be represented
as a function of r, i.e., wΩ(r). To this end, some recent
studies [4, 16, 52] defined baseline values from a heuris-
tic perspective, e.g. simply using mean/zero baseline val-
ues [16, 67]. However, it still remains an open problem to
define optimal baseline values.
Thus, we further boost the sparsity of causal patterns
by learning the optimal baseline values that enhance the
conciseness of the causal graph. However, it is difficult to
learn optimal baseline values by directly optimizing Eq. (5).
To this end, we relax the optimization problem in Eq. (5) (L0
regression) as a Lasso regression (L1 regression) as follows.
4
